% block4_psor.tex
%
% Phase 3, Block 4: American put pricing via PSOR.
%
% Self-contained writeup. Compiles standalone with pdflatex.

\documentclass[11pt,a4paper]{article}

\usepackage[T1]{fontenc}
\usepackage[utf8]{inputenc}
\usepackage{amsmath, amssymb, amsthm}
\usepackage{geometry}
\usepackage{hyperref}
\usepackage{enumitem}

\geometry{margin=1in}

\newtheorem{proposition}{Proposition}[section]
\newtheorem{theorem}[proposition]{Theorem}
\newtheorem{remark}[proposition]{Remark}
\newtheorem{definition}[proposition]{Definition}

\title{Phase 3, Block 4:\\Pricing the American put via PSOR}
\author{Quant Finance Portfolio}
\date{}

\begin{document}
\maketitle

\section{Setting and motivation}

The American put can be exercised at any time $t \in [0, T]$, not
only at maturity. This single change relative to the European case
transforms the pricing problem qualitatively: the price now depends
on a free \emph{exercise boundary} $S_f(t)$ that separates a
\emph{continuation region} (where the option is held) from an
\emph{exercise region} (where it is exercised immediately). The
boundary itself is unknown a priori --- it is part of the solution.

The deliverable of this block is a numerical method that handles the
free boundary cleanly, by reformulating the problem as a
\emph{Linear Complementarity Problem} (LCP) on a fixed domain and
solving it iteratively at every time step using \emph{Projected
Successive Over-Relaxation} (PSOR). The block ships two modules:
\begin{itemize}[leftmargin=2em]
\item \texttt{quantlib.psor}: a generic PSOR solver for tridiagonal
LCPs (analogous to \texttt{quantlib.thomas} of Block~2).
\item \texttt{quantlib.cn\_american}: the user-facing American put
pricer, built on the $\theta$-scheme infrastructure of Block~3 with
PSOR replacing Thomas in the per-step solve.
\end{itemize}

% =====================================================================
\section{The American put as an optimal stopping problem}
% =====================================================================

Under the risk-neutral measure $\mathbb{Q}$, the price of the
American put is
\begin{equation}
V(S_0, 0) = \sup_{\tau \in \mathcal{T}_{[0,T]}}
\mathbb{E}^{\mathbb{Q}}\!\left[\,e^{-r\tau}(K - S_\tau)^+ \,\big|\,
S_0\,\right],
\label{eq:opt-stop}
\end{equation}
where $\mathcal{T}_{[0,T]}$ is the set of stopping times taking
values in $[0, T]$. Equation \eqref{eq:opt-stop} states that the
holder picks the stopping rule $\tau$ that maximises the discounted
expected payoff.

The optimal $\tau$ partitions the state space $(S, t)$ into two
regions:
\begin{itemize}[leftmargin=2em]
\item \textbf{Continuation region} $\mathcal{C} = \{(S, t) : V(S, t) >
(K - S)^+\}$: holding the option is strictly better than exercising
now.
\item \textbf{Exercise region} $\mathcal{E} = \{(S, t) : V(S, t) =
(K - S)^+\}$: it is optimal (or at least no worse) to exercise.
\end{itemize}
The boundary between the two is the curve
$S = S_f(t)$, the \emph{exercise boundary}, also called the
\emph{free boundary} because its location is part of the solution
rather than known a priori. For a put, $S_f(t)$ is increasing in
$t$ with $S_f(T^-) = K$.

% =====================================================================
\section{The classical free-boundary PDE formulation}
% =====================================================================

In the continuation region, the standard arbitrage argument shows
that $V$ satisfies the Black--Scholes PDE:
\begin{equation}
\frac{\partial V}{\partial t} + \tfrac{1}{2}\sigma^2 S^2
\frac{\partial^2 V}{\partial S^2} + r S \frac{\partial V}{\partial S}
- r V = 0, \qquad (S, t) \in \mathcal{C}.
\label{eq:bs-pde-cont}
\end{equation}
In the exercise region, $V$ equals the intrinsic value:
\begin{equation}
V(S, t) = K - S, \qquad (S, t) \in \mathcal{E}.
\label{eq:exercise-region}
\end{equation}
At the free boundary $S = S_f(t)$, two matching conditions hold:
\begin{align}
V(S_f(t), t) &= K - S_f(t), \label{eq:value-matching} \\
\frac{\partial V}{\partial S}(S_f(t), t) &= -1.
\label{eq:smooth-pasting}
\end{align}
The first is value matching: $V$ is continuous across the boundary.
The second is the \emph{smooth pasting} condition: $V_S$ is also
continuous. Two conditions are needed because the boundary itself
is unknown; smooth pasting is the optimality condition that
determines $S_f(t)$.

\paragraph{Why this formulation is hard numerically.} The PDE
\eqref{eq:bs-pde-cont} holds on the time-varying domain
$\mathcal{C}$, whose boundary moves with $t$. The finite-difference
methods of Blocks 1--3 work on fixed domains; adapting the spatial
mesh to track $S_f(t)$ would require remeshing at every time step,
which is expensive and error-prone. We need a reformulation on a
fixed domain.

% =====================================================================
\section{The variational reformulation as an LCP}
% =====================================================================

Switching to log-coordinates $x = \ln(S/K)$ and $\tau = T - t$ and
denoting the spatial operator
$\mathcal{L}V = \frac{\sigma^2}{2}V_{xx} + \mu V_x - rV$ as in
Block~3, we reformulate the American put as the system
\begin{equation}
\boxed{\;\;\begin{aligned}
V_\tau - \mathcal{L}V \;&\geq\; 0 \\
V - g \;&\geq\; 0 \\
(V_\tau - \mathcal{L}V)(V - g) \;&=\; 0
\end{aligned}\;\;}
\qquad
\text{on } x \in (x_{\min}, x_{\max}),\; \tau \in (0, T],
\label{eq:lcp-cont}
\end{equation}
where $g(x) = K(1 - e^x)^+$ is the put payoff function (the
\emph{obstacle}). Initial condition: $V(x, 0) = g(x)$. Boundary
conditions: $V(x_{\min}, \tau) = K - K e^{x_{\min}}$ (deep ITM:
exercise immediately, no discount) and
$V(x_{\max}, \tau) = 0$ (deep OTM: worthless).

\paragraph{Why each inequality is correct.} Take them in turn.

\begin{enumerate}[leftmargin=2em]
\item $V_\tau - \mathcal{L}V \geq 0$. In $\mathcal{C}$ this is an
equality (the BS PDE itself). In $\mathcal{E}$ we have $V \equiv g$,
so $V_\tau = 0$, and a direct calculation
$\mathcal{L}g = -rK \leq 0$ for $x < 0$ (the put-ITM region) makes
$V_\tau - \mathcal{L}V = -\mathcal{L}g \geq 0$.

\item $V - g \geq 0$. The price never falls below the immediate
payoff: if it did, exercising would dominate holding, an arbitrage.

\item $(V_\tau - \mathcal{L}V)(V - g) = 0$ \emph{(complementarity)}.
At every $(x, \tau)$, exactly one of the two preceding inequalities
is strict: in $\mathcal{C}$ the second is strict ($V > g$) and the
first is an equality (BS PDE); in $\mathcal{E}$ the first is strict
($V_\tau - \mathcal{L}V > 0$) and the second is an equality
($V = g$).
\end{enumerate}

\paragraph{The structural payoff.} The system \eqref{eq:lcp-cont}
makes no reference to the free boundary $S_f$. The boundary emerges
a posteriori as the locus where one inequality changes from strict
to equality. Numerically, this is enormous: we can work on the
fixed log-truncated domain of Block~0, with no need to track a
moving boundary, and the complementarity condition automatically
selects which nodes are in which region.

\paragraph{Equivalence.} The LCP \eqref{eq:lcp-cont} is equivalent
to the classical PDE+free-boundary formulation
\eqref{eq:bs-pde-cont}--\eqref{eq:smooth-pasting} for solutions
$V \in C^1$. The smooth pasting condition is recovered automatically:
from $V \geq g$ everywhere and $V = g$ on $\mathcal{E}$, taking the
spatial derivative at the boundary gives $V_S \geq g_S = -1$ from
the continuation side and $V_S = -1$ on the exercise side; combined
with $V \in C^1$, this forces $V_S = -1$ at the boundary. Reference:
Wilmott--Howison--Dewynne~\cite{WHD1995}, Chapter~7.

% =====================================================================
\section{Discretisation: from continuous to discrete LCP}
% =====================================================================

Apply the $\theta$-scheme to the operator $V_\tau - \mathcal{L}V$
exactly as in Block~3. The discrete equality $A V^{n+1} = B V^n +
\mathbf{b}_{\mathrm{bnd}}$ (with $A$, $B$, $\mathbf{b}_{\mathrm{bnd}}$ from
Block~3) becomes the discrete LCP
\begin{equation}
\boxed{\;\;\begin{aligned}
A V^{n+1} \;&\geq\; B V^n + \mathbf{b}_{\mathrm{bnd}} \\
V^{n+1} \;&\geq\; \mathbf{g} \\
\bigl(A V^{n+1} - B V^n - \mathbf{b}_{\mathrm{bnd}}\bigr) \cdot
\bigl(V^{n+1} - \mathbf{g}\bigr) \;&=\; 0
\end{aligned}\;\;}
\label{eq:lcp-discrete}
\end{equation}
componentwise, where $\mathbf{g} \in \mathbb{R}^{N-1}$ is the payoff
sampled on the interior grid: $g_j = K(1 - e^{x_j})^+$.

Each time step now requires solving the LCP \eqref{eq:lcp-discrete}
instead of a linear system. Thomas does not apply: it solves
$Ax = b$, not $Ax \geq b$ with $x \geq c$. We need a new solver.

% =====================================================================
\section{SOR: the iterative method behind PSOR}
% =====================================================================

For the linear system $Ax = b$ with $A$ tridiagonal and diagonal
dominant, an alternative to Thomas is the iterative
\emph{Successive Over-Relaxation} (SOR) method. SOR produces a
sequence $x^{(0)}, x^{(1)}, x^{(2)}, \dots$ that converges to the
exact solution. Although asymptotically more expensive than Thomas
($O(N \log\epsilon^{-1})$ per solve vs $O(N)$), SOR has a structural
advantage: it generalises naturally to LCPs.

\paragraph{Definition.} SOR updates one component at a time. For a
tridiagonal system $a_{i,i-1} x_{i-1} + a_{ii} x_i + a_{i,i+1}
x_{i+1} = b_i$, SOR uses
\begin{equation}
x_i^{(k+1)} = (1 - \omega)\,x_i^{(k)} \,+\,
\omega \cdot \frac{1}{a_{ii}}\Bigl(\, b_i \,-\, a_{i,i-1}\,
x_{i-1}^{(k+1)} \,-\, a_{i,i+1}\,x_{i+1}^{(k)} \,\Bigr)
\label{eq:sor}
\end{equation}
for $i = 1, \dots, n$, with $\omega \in (0, 2)$ the
\emph{relaxation parameter}. At $\omega = 1$, equation \eqref{eq:sor}
is plain Gauss--Seidel; at $\omega > 1$ it over-relaxes (extrapolates
beyond the GS update); at $\omega < 1$ it under-relaxes
(interpolates between old and new). Note that \eqref{eq:sor} uses
$x_{i-1}^{(k+1)}$ (the just-updated value) and $x_{i+1}^{(k)}$ (the
old value) --- this in-place update is what makes SOR efficient.

\paragraph{Convergence.} A classical result (Kahan; Ostrowski--Reich):

\begin{theorem}[SOR convergence]
If $A$ is symmetric positive definite, SOR converges for every
$\omega \in (0, 2)$ from any initial guess $x^{(0)}$. For a
tridiagonal $A$ with constant coefficients (our case), the
convergence rate is geometric with optimal relaxation parameter
\begin{equation}
\omega^* = \frac{2}{1 + \sqrt{1 - \rho_J^2}},
\label{eq:omega-star}
\end{equation}
where $\rho_J$ is the spectral radius of the Jacobi iteration matrix
$D^{-1}(L + U)$.
\end{theorem}

For the matrices arising from BTCS or CN, $\rho_J \approx 1 - O(1/N)$,
giving $\omega^* \approx 2 - O(1/\sqrt{N})$. For typical project
parameters, $\omega^* \in [1.2, 1.5]$.

\paragraph{Asymmetry.} The CN matrix $A$ is not symmetric (the
convective term $\mu V_x$ uses centred differences, breaking
symmetry). The Ostrowski--Reich theorem does not apply directly.
However, $A$ remains diagonally dominant, and a generalisation
(Varga~\cite{Varga2000}, Theorem~3.5) covers diagonally dominant
M-matrices: SOR converges for $\omega \in (0, 2)$, with
$\omega^*$ moderately shifted from \eqref{eq:omega-star}. In
practice, for our parameter range with grid Pe'clet $\leq 0.25$,
$\omega \approx 1.2$ is near optimal.

% =====================================================================
\section{PSOR: SOR with projection}
% =====================================================================

PSOR is the generalisation of SOR to linear complementarity problems.
The modification is a single line: after each SOR component update,
project onto the constraint $x_i \geq g_i$:
\begin{equation}
\boxed{\;\;
x_i^{(k+1)} = \max\!\left\{\,g_i,\;\;
(1 - \omega)\,x_i^{(k)}
+ \omega \cdot \frac{1}{a_{ii}}\Bigl(b_i
- a_{i,i-1}\,x_{i-1}^{(k+1)}
- a_{i,i+1}\,x_{i+1}^{(k)}\Bigr)\,\right\}.
\;\;}
\label{eq:psor}
\end{equation}
That is: compute the SOR candidate, and if it falls below the
obstacle $g_i$, raise it to $g_i$. If it lies above, leave it.

\paragraph{Variational interpretation.} The LCP
\eqref{eq:lcp-discrete} is equivalent (when $A$ is symmetric positive
definite) to the constrained minimisation
\[
\min_{x \geq \mathbf{g}}\;\; \tfrac{1}{2} x^T A x - x^T b,
\]
where $b = B V^n + \mathbf{b}_{\mathrm{bnd}}$. PSOR is the
componentwise application of SOR to the KKT conditions of this
minimisation, and the projection $\max\{g_i, \cdot\}$ is exactly the
projection onto the feasible set $\{x : x \geq \mathbf{g}\}$. This
construction is due to Cryer~\cite{Cryer1971}.

\begin{theorem}[Cryer 1971]
If $A \in \mathbb{R}^{n \times n}$ is symmetric positive definite
and $\omega \in (0, 2)$, the PSOR iteration \eqref{eq:psor} converges
to the unique solution of the LCP $A x \geq b$, $x \geq \mathbf{g}$,
$(Ax - b)(x - \mathbf{g}) = 0$ from any initial guess $x^{(0)} \geq
\mathbf{g}$.
\end{theorem}

\paragraph{For non-symmetric $A$.} The CN matrix is not symmetric,
but is a strictly diagonally dominant M-matrix in the project's
parameter range. PSOR convergence in this setting follows from
results of Mangasarian and others; see Cottle--Pang--Stone
\cite{CPS2009}, \S5.3 and \S5.6 for the rigorous treatment. The
qualitative picture is that for moderate Pe'clet number, asymmetry
slows convergence by a constant factor but does not destroy it.

\paragraph{Initialisation.} A natural choice is $x^{(0)} = \mathbf{g}$
(the obstacle itself). It is feasible by construction, and is exact
on the eventual exercise region. An alternative is warm-starting from
the previous time step, $x^{(0)} = V^n$. Since $V^{n+1} - V^n =
O(\Delta\tau)$, the warm start saves roughly half the iterations
in practice.

% =====================================================================
\section{Stopping criterion}
% =====================================================================

The iteration must be terminated at finite $k$. The standard test
is on the increment:
\begin{equation}
\bigl\|x^{(k+1)} - x^{(k)}\bigr\|_\infty
\,<\,
\texttt{tol}_{\mathrm{abs}}
+ \texttt{tol}_{\mathrm{rel}} \cdot \bigl\|x^{(k)}\bigr\|_\infty,
\label{eq:stopping}
\end{equation}
with typical values $\texttt{tol}_{\mathrm{abs}} = 10^{-8}$ and
$\texttt{tol}_{\mathrm{rel}} = 10^{-7}$. The absolute component
protects against deep-OTM puts where prices are small (cents); the
relative component protects against deep-ITM puts where prices are
large (close to $K$). The mixed criterion is robust across the entire
strike range.

\paragraph{Why not the residual.} The residual criterion
$\|Ax - b\| < \texttt{tol}$, standard for linear systems, does not
apply to the LCP: the residual is non-zero in the exercise region
(by the first inequality) even at the optimum. The increment
criterion is agnostic to the problem type and works for both linear
systems and LCPs.

\paragraph{Maximum iteration safeguard.} If $\omega$ is mis-tuned or
the matrix is poorly conditioned, the iteration may converge slowly
or not at all. A hard limit \texttt{max\_iter} (typically $10^4$)
prevents an infinite loop. Reaching the limit raises a
\texttt{RuntimeError} (or \texttt{std::runtime\_error} in C++),
making a silent failure visible to the caller.

% =====================================================================
\section{Computational cost}
% =====================================================================

Per time step:
\begin{itemize}[leftmargin=2em]
\item Each PSOR sweep is $O(N)$ operations.
\item Number of sweeps to reach tolerance $\epsilon$:
$K \sim -\log\epsilon / \log(1/\rho(\omega^*))$.
For $\rho(\omega^*) \approx 1 - C/N$ (typical for our matrices)
and $\epsilon = 10^{-7}$, $K = O(N)$.
\item Cost per time step: $O(N^2)$.
\end{itemize}
Total cost over $M$ time steps: $O(M \cdot N^2)$. Taking $M \sim N$
(matched to spatial accuracy as in CN), this is $O(N^3)$. For
$N = 200$, $\epsilon = 10^{-7}$, the constant ratio of CN-PSOR to CN-Thomas is
roughly $K \sim 200$ --- two orders of magnitude more work per step
than the European pricer, the price of LCP generality.

\paragraph{Comparison with other American pricers.}
\begin{center}
\begin{tabular}{lll}
\hline
Method & Cost & Accuracy \\
\hline
CRR binomial (Phase 2)        & $O(N^2)$ & $O(1/N)$, $O(1/N^2)$ Richardson \\
LSM (Phase 2)                 & $O(N_{\mathrm{paths}} N_{\mathrm{steps}})$ & MC, $O(1/\sqrt{N_{\mathrm{paths}}})$ \\
\textbf{CN-PSOR (this block)} & $O(N^3 \log\epsilon^{-1})$ & $O(\Delta x^2 + \Delta\tau^2)$ \\
\hline
\end{tabular}
\end{center}
CN-PSOR is the most accurate but most expensive at moderate $N$.
LSM scales best to high dimensions; CRR is the simplest and most
robust. Block~6 makes the comparison numerical.

% =====================================================================
\section{Implementation outline}
% =====================================================================

Two new modules:

\paragraph{\texttt{quantlib.psor} / \texttt{quant::pde::psor}.}
Generic LCP solver. The signature mirrors Thomas:
\begin{verbatim}
def psor_solve(sub, diag, sup, rhs, obstacle, *,
               omega=1.2, tol_abs=1e-8, tol_rel=1e-7,
               max_iter=10000, x0=None) -> (x, n_iter)
\end{verbatim}
Returns both the solution and the number of iterations performed
(the iteration count is informative diagnostic data; the lesson
section below uses it to demonstrate $\omega$ optimality).

\paragraph{\texttt{quantlib.cn\_american} / \texttt{quant::pde::cn\_american}.}
American put pricer. Reuses the $\theta$-scheme grid and boundary
construction of Block~3, replacing the per-step Thomas solve with a
PSOR solve. Two notable differences from \texttt{cn\_european\_put}:
\begin{enumerate}[leftmargin=2em]
\item The lower boundary condition is $V(x_{\min}, \tau) = K
- K e^{x_{\min}}$ (no discount, since exercise is immediate).
\item No Rannacher smoothing. Justification: the Rannacher protocol
addresses the kink in the European payoff, which produces
high-frequency oscillations that CN cannot damp. For the American
put the oscillations are the same in nature, but they are
\emph{absorbed by the projection step}: nodes near the strike fall
in the exercise region $\mathcal{E}$ where $V = g$ exactly, so the
oscillation is clipped at $g$ before it can propagate. Empirically,
CN-PSOR for the American put has $O(\Delta\tau^2)$ convergence
without Rannacher; the projection does the smoothing implicitly.
\end{enumerate}

% =====================================================================
\section{Validation strategy}
% =====================================================================

Seven tests certify the new pricer:

\begin{enumerate}[leftmargin=2em]
\item \textbf{Cross-validation against CRR binomial} (from Phase 2):
the gold-standard reference for American puts. Tolerance comparable
to the discretisation error.

\item \textbf{Cross-validation against LSM} (from Phase 2): a
secondary check. Tolerance loosened to account for Monte Carlo
variance.

\item \textbf{Quadratic convergence in $\Delta x$}: doubling $N$
should reduce the error by approximately 4. Demonstrates that the
$\theta = 1/2$ scheme retains its second-order spatial accuracy under
the projection.

\item \textbf{American $\geq$ European}: at every spot, the American
put price must be at least the European put price. Strict in the
exercise region, equality possible deep in the continuation region.

\item \textbf{$\omega$-sweep convergence}: at fixed problem,
$\omega \in \{0.5, 1.0, 1.2, 1.5, 1.9\}$ produces different
iteration counts. The minimum should occur near $\omega \approx 1.2$,
confirming the theoretical $\omega^*$ prediction.

\item \textbf{$\omega$ outside $(0, 2)$ raises}: $\omega = 0$ or
$\omega = 2$ produces no-convergence within
\texttt{max\_iter}, raising a runtime error.

\item \textbf{Free boundary monotonicity}: the recovered exercise
boundary $S_f(\tau)$ (the largest $S$ such that $V = g$ at each
$\tau$) is monotonically increasing in $\tau$, with $S_f(\tau \to T)
\to K$. A qualitative shape test that catches gross errors.
\end{enumerate}

% =====================================================================
\section*{References}
% =====================================================================

\begin{thebibliography}{9}

\bibitem{WHD1995}
P. Wilmott, S.~Howison and J. Dewynne.
\emph{The Mathematics of Financial Derivatives.}
CUP, 1995. Chapter 7 on American options as free-boundary problems;
self-contained derivation of smooth pasting from optimality.

\bibitem{Wilmott1998}
P. Wilmott.
\emph{Paul Wilmott on Quantitative Finance.}
Wiley, 1998 (2nd ed. 2006). Chapters 8 and 35: practical PSOR
implementation with code.

\bibitem{TavellaRandall2000}
D. Tavella and C. Randall.
\emph{Pricing Financial Instruments: The Finite Difference Method.}
Wiley, 2000. Chapter 5 on American options; explicit treatment of
PSOR with code.

\bibitem{Seydel2017}
R. Seydel.
\emph{Tools for Computational Finance.}
6th ed., Springer, 2017. \S4.5 on American options; concise.

\bibitem{Cryer1971}
C.~W. Cryer.
\emph{The solution of a quadratic programming problem using
systematic overrelaxation.}
SIAM Journal on Control 9, 385--392, 1971. Original PSOR paper:
proof of convergence for symmetric positive-definite LCPs.

\bibitem{CPS2009}
R.~W. Cottle, J.-S. Pang and R.~E. Stone.
\emph{The Linear Complementarity Problem.}
SIAM Classics in Applied Mathematics, 2009 (orig. Academic Press,
1992). The rigorous reference for LCP theory. \S5.2--5.3 cover
PSOR convergence in detail.

\bibitem{Glowinski2008}
R. Glowinski.
\emph{Numerical Methods for Nonlinear Variational Problems.}
Springer, 2008 (orig. Springer 1984). The variational treatment of
obstacle problems.

\bibitem{Varga2000}
R.~S. Varga.
\emph{Matrix Iterative Analysis.}
Revised 2nd ed., Springer, 2000. Chapters 3--4 on SOR convergence
theory and optimal $\omega$.

\bibitem{Saad2003}
Y. Saad.
\emph{Iterative Methods for Sparse Linear Systems.}
2nd ed., SIAM, 2003. Chapter 4 on relaxation methods, modern
treatment.

\end{thebibliography}

\end{document}
